\renewcommand{\theposttitle}{The Cycle of Growth and Stoic Philosophy}
\renewcommand{\thepostdate}{2022-07-15}
\vspace*{\fill}
\begin{center}
\includegraphics[width=\textwidth, height=0.55\textheight, keepaspectratio]{C:/Users/James/Documents/GitHub/VisualIdeasBook/data/images/f635d99b631a1ba5ff075e5cce550266.png}
\end{center}
\vspace{2mm}
\begin{center}
\begin{minipage}{0.92\textwidth}
\small
This image made it to the Interesting slot because it's the first image I posted where I received a ton of feedback on the analogy itself (which I love actually because it forces me to improve my thinking and iterate the drawing).

For me, the key idea is when you're working on something alone, the low-hanging fruit is the easiest to grab.

But if you want a sweeter reward, getting there in a team is easier, even if the fruit (🥁) of that labor may need to be shared.

Honestly though...working together with a community toward a goal is more fun and better for the spirit.

I imagine the characters in this picture cracking up as they try not to topple over one another. 😂

\textbf{Artist's Note:} Iteration for Clarity• Include all three of the characters grabbing low-hanging fruit with three golden apples at the top. • When they're grabbing the fruit on the right, "Bluey" will have a basket full of 2 golden apples as he holds the third.
\end{minipage}
\end{center}
\vspace*{\fill}